% ===========================================
At the beginning, my personal motivation was to build a list of references, each with a summary of the results therein related to Wilf's conjecture. This would have helped me by not having to dive into a collection of papers each time I needed a result. Then I thought that making the list public could also be a contribution to Wilf's conjecture. This process ended up in the writing of this paper, which is in some sense \emph{yet another survey}. Another, because most papers fully dedicated to the conjecture provide good literature reviews. Although not aiming to be complete, these could be taken as surveys.

As it frequently happens with easy to state combinatorial problems, while working on them one thinks that a solution is at reach. This is certainly the case of the problem posed by Wilf, but nevertheless no one has found the aimed solution so far. Tacking into account the number of published papers on the theme, one can infer that much time has globally been dedicated to the problem. This may lead people to classify the problem in the category of dangerous problems, in the sense that one risks to spend too much time struggling with it and have to give up without getting a solution. Fortunately, partial results may be of some interest.

The plan of the paper follows.

This introductory section contains most of the terminology and notation to be used along the paper. 
There are not many differences to what is commonly used. 
This section contains also what I consider a convenient way to visualize numerical semigroups. Although almost all further images appear only in the last section, we provided sufficient information to produce images of semigroups appearing in the remaining parts of the text.

Some problems posed by Wilf are described in the second section, which can be seen as a kind of motivation for the paper.

The third section is the real survey. It contains a large introductory part and then the statements of results grouped into several subsections.

In the final section we introduce a notion of \emph{quasi generalization} (roughly speaking, a set quasi generalizes another if it contains all its elements, except possibly a finite number of them). It allows to draw a lattice involving some important properties that give rise to semigroups satisfying Wilf's conjecture. 
% ===========================================

